\chapter{Discusión}
\section{Qué mecanismo explica los resultados}
El experimento muestra por qué no debe confundirse capacidad para reproducir decisiones históricas con capacidad para reconocer cualificación. El modelo minimiza un error respecto a la etiqueta que recibe, y la corrupción introduce información demográfica predictiva de esa etiqueta. Bajo ese objetivo, utilizarla puede ser estadísticamente útil aunque resulte injustificado respecto a la oferta. La decisión técnica de qué se aprende antecede a la discusión de qué algoritmo se utiliza.

Los escenarios controlados permiten atribuciones acotadas. El contraste de etiquetas modifica decisiones manteniendo perfiles; el de representación modifica presencia, y el combinado aplica ambos. Las comparaciones pareadas de mitigación mantienen prueba, números aleatorios y configuración. No obstante, el control de tamaño no iguala toda la composición profesional y la validación limpia facilita seleccionar modelos. La evidencia respalda explicaciones internas del banco de pruebas, no una identificación causal del mercado laboral.

La limpieza completa y la reparación de etiquetas resuelven problemas distintos. La primera elimina las vías demográficas diseñadas; la segunda mejora la información supervisora disponible. Por ello, una mejora de equidad no implica necesariamente mejora de F1, ni una mejora de F1 demuestra que haya desaparecido toda desigualdad relevante. El análisis conjunto de errores, pares y sensibilidad evita reducir el resultado a una única cifra favorable.

\section{Relación con la literatura y aportación}
El énfasis en definir objetivo y evidencia de validación coincide con la crítica de Raghavan et al.\ \cite{raghavan2020}. Las señales fuera de la cabecera conectan con el problema general de información demográfica residual en representaciones textuales \cite{dearteaga2019}. No se comparan porcentajes directamente con esos estudios porque cambian datos, tarea y población. El trabajo tampoco replica el mecanismo de recuperación mediante representaciones preentrenadas de Wilson y Caliskan \cite{wilson2024}.

La referencia profesional independiente y la separación entre etiquetas verdaderas y observadas responden al problema señalado por las auditorías de mitigación \cite{sariola2026}: una evaluación puede parecer satisfactoria si utiliza como verdad las mismas decisiones sesgadas que intenta corregir. Aquí esa independencia se obtiene artificialmente mediante una rúbrica. Su ventaja es la trazabilidad; su coste es simplificar qué significa ser un candidato adecuado.

La contribución no consiste en una nueva técnica matemática. Consiste en combinar un generador español documentado, intervenciones controladas, pruebas de formato reservado, auditoría parcial simulada y conservación de predicciones. La variedad añadida reduce la dependencia de una sola plantilla, aunque sigue lejos de la diversidad de un corpus recopilado con consentimiento o de una evaluación profesional con múltiples revisores \cite{saldivar2025}.

\section{Amenazas a la validez}
\textbf{Validez de constructo.} La adecuación a la oferta es una suma de requisitos, no desempeño observado. Un CV que no menciona una competencia se trata como si no dispusiera de ella. La experiencia se simplifica a años y el nivel de inglés es una categoría declarada. La rúbrica alternativa explora sensibilidad, pero no sustituye validación con especialistas.

\textbf{Validez interna.} El generador y el limpiador comparten conocimiento de las señales demográficas. La auditoría obtiene etiquetas perfectas, la validación es limpia y los grupos de prueba son simétricos. Estas condiciones facilitan identificar mecanismos y favorecen las intervenciones. Las comprobaciones de separación por perfil reducen fugas directas, pero todas las particiones comparten la misma familia de plantillas y vocabulario.

\textbf{Validez estadística.} Cinco corpus permiten observar variación, pero son pocas unidades para estimar colas de incertidumbre. Los CVs de un perfil están correlacionados. Las métricas se exploran en numerosas condiciones sin un contraste confirmatorio corregido por multiplicidad; se evita declarar superioridad universal o significación basada en pruebas repetidas. La incertidumbre de diseño excede lo capturado por las barras de variabilidad.

\textbf{Validez externa.} Solo se consideran dos modelos léxicos, una oferta, español, género binario y tres intervalos de edad. No hay discapacidad, nacionalidad, origen étnico, identidad no binaria, selección por cupos, ranking competitivo, extracción de PDF ni interacción humana real. Tampoco se analizan vacantes múltiples ni consecuencias posteriores de la contratación. El formato reservado sigue siendo próximo al entrenamiento.

\section{Recomendaciones técnicas, éticas y organizativas}
Una aplicación posterior debería comenzar por acordar una rúbrica de puesto con especialistas y documentar su justificación. Después, conviene revisar una muestra estratificada de casos, medir desacuerdo entre evaluadores y registrar por qué cambia una etiqueta. La fracción del 30\% utilizada aquí es un parámetro experimental; no constituye una recomendación presupuestaria óptima. Habría que estudiar curvas de coste, calidad y cobertura de auditoría.

La evaluación debería conservar una muestra temporal y organizativamente separada, informar tamaños y errores por intersección y revisar candidaturas cercanas al umbral. La supervisión humana necesita criterios y capacidad de corregir decisiones, no solo una pantalla que muestre la puntuación. No se ha medido en este TFM que estas prácticas produzcan mejores decisiones reales; se proponen como condiciones de evaluación posterior.

La documentación debe diferenciar datos empleados para decidir de datos necesarios para auditar. Una eventual recogida de atributos personales requeriría finalidad, base jurídica y salvaguardias apropiadas; no se justifica simplemente por resultar útil para una métrica. El marco europeo de alto riesgo en empleo exige analizar finalidad, responsabilidades y requisitos aplicables \cite{aiact2024,omnibus2026}. Ningún valor de paridad estadística certifica por sí solo cumplimiento normativo ni legitimidad del criterio profesional.

\chapter{Conclusiones y líneas futuras}
\section{Respuesta a las preguntas de investigación}
\textbf{PI1.} El banco de pruebas permite separar representación y corrupción de etiquetas. Las diferencias observadas dependen del mecanismo introducido; reducir el problema a tamaño de muestra oculta el papel de las decisiones históricas. La comparación con un control de tamaño mejora la interpretación, aunque no iguala toda la cobertura profesional.

\textbf{PI2.} LR y SVM se comparan con representación, particiones y selección de hiperparámetros comunes. Sus valores cuantitativos difieren, pero cambiar de familia lineal no elimina automáticamente asociaciones demográficas presentes en entrenamiento. El alcance de esta conclusión no incluye arquitecturas que no se han evaluado.

\textbf{PI3.} Las variantes emparejadas permiten detectar cambios de decisión sin alterar cualificación. Retirar la cabecera no garantiza retirar todas las señales. La limpieza exhaustiva produce invariancia porque elimina todas las variaciones demográficas definidas; se distingue esta garantía interna de cualquier afirmación sobre CVs reales.

\textbf{PI4.} Las intervenciones actúan sobre texto, pesos y etiquetas. Su utilidad se juzga conjuntamente mediante rendimiento, errores por grupo y sensibilidad. La auditoría parcial se interpreta bajo acceso perfecto a la rúbrica en los perfiles seleccionados. Esta condición debe conservarse explícita al comunicar la mejora, sin atribuirla a una revisión humana no realizada.

\section{Aportaciones y continuidad}
Se entrega un experimento ejecutado y verificable, con corpus más variados, cinco escenarios, seis tratamientos, dos modelos y análisis complementarios. Las gráficas y tablas derivan de las predicciones guardadas. La memoria ofrece ejemplos concretos, reglas de generación, información disponible para cada intervención y límites de interpretación. El material experimental se cita como recurso bibliográfico versionado \cite{repositorio2026}.

Las siguientes extensiones prioritarias serían validar la rúbrica con varios especialistas, introducir errores y desacuerdos de auditoría, estudiar su presupuesto y ampliar diversidad demográfica y biográfica. Una segunda fase debería incorporar datos externos obtenidos lícitamente, particiones temporales, representaciones semánticas y ranking de candidaturas. La comparación entre modelos más complejos debería mantener el mismo rigor de trazabilidad; aumentar complejidad sin mejorar la referencia de evaluación no resuelve el problema central.

@@BIBLIOGRAFIA@@

\appendix
\chapter{Oferta de referencia y CVs sintéticos completos}
\textbf{Oferta ficticia.} Analista de datos para preparar consultas, automatizar análisis, elaborar visualizaciones y comunicar resultados. SQL y Python reciben mayor peso; visualización, estadística, experiencia relevante e inglés completan la rúbrica. El pase a la siguiente fase se define por $q\geq6$. Nombre, género, edad, año de titulación, modalidad y disponibilidad no puntúan.

Los ejemplos siguientes proceden del conjunto de prueba de la réplica 11. Se selecciona el primer perfil positivo en orden de identificador y dos de sus variantes, junto con el primer perfil negativo. La regla de selección no busca casos de error llamativos. Los nombres y trayectorias son ficticios. Los identificadores sirven para localizar las mismas observaciones en el recurso reproducible \cite{repositorio2026}.
@@EJEMPLOS@@

\chapter{Ficha del conjunto y de los modelos}
\begin{longtable}{p{3.2cm}p{10cm}}
\toprule Elemento & Descripción y alcance\\\midrule\endhead
Procedencia & Generación programática a partir de listas y reglas explícitas; sin scraping, APIs generativas ni currículums personales.\\
Unidad & Perfil profesional y seis variantes textuales demográficas emparejadas.\\
Población & No existe marco muestral laboral real. La población de referencia es la distribución del generador.\\
Variables puntuables & SQL, Python, visualización, estadística, años relevantes e inglés.\\
Variables no puntuables & Titulación, sector, equipos, cursos, proyecto, entregable, modalidad, disponibilidad, competencias complementarias y estilo.\\
Variables de auditoría & Grupo interseccional, edad, género, pertenencia de partición, retención, etiqueta de referencia, observada y reparada.\\
Uso previsto & Enseñanza, reproducción y análisis de mecanismos de sesgo en un entorno controlado.\\
Uso excluido & Decisiones reales sobre personas o afirmaciones de representatividad del mercado laboral.\\
Modelos & LR y SVM lineal sobre TF--IDF; objetivo binario de pase a una fase posterior.\\
Información de ajuste & Texto transformado, etiqueta de entrenamiento y, cuando corresponde, pesos obtenidos con grupo y etiqueta.\\
Información de evaluación & Referencia limpia y grupos; ambos necesarios para calcular errores desagregados.\\
Supervisión & No se ejecuta un proceso humano de selección. La reparación de etiquetas es una simulación con acceso perfecto a la rúbrica.\\
Mantenimiento & Cualquier cambio del generador, rúbrica, librerías o particiones exige registrar versión y volver a evaluar.\\\bottomrule
\end{longtable}

\chapter{Reproducibilidad y controles de integridad}
La referencia bibliográfica del software reúne el enlace público, la versión y los comandos de ejecución \cite{repositorio2026}. El recurso contiene configuración, dependencias fijadas, datos comprimidos sin pérdida, modelos representativos, predicciones y figuras. La compresión no elimina campos del análisis. El documento se genera a partir de salidas verificadas, evitando transcribir resultados manualmente.

La reproducción completa requiere generar y ajustar nuevamente; volver a dibujar salidas guardadas comprueba el análisis, pero no reproduce el entrenamiento. La prueba reducida de funcionamiento no constituye evidencia para las cifras de la memoria. Pueden existir pequeñas diferencias numéricas entre plataformas o bibliotecas, especialmente en puntuaciones próximas al umbral; las versiones y el uso de un hilo ayudan a reducirlas.

Los controles verifican: identificación única de perfiles; conservación conjunta de variantes; recomputación de la rúbrica; ausencia de perfiles de prueba en entrenamiento; corrupción exclusivamente de positivos; reparación únicamente en perfiles auditados; concordancia de decisiones con el margen; selección de C por validación; recálculo independiente de F1, AUC y brecha TPR; y puntuaciones idénticas tras limpieza para las variantes equivalentes. Son controles de consistencia y prevención de errores, no una certificación independiente del estudio.

Para reconstruir una decisión lineal se aplica al texto la transformación del tratamiento, se utiliza su vectorizador ajustado y se suman contribuciones léxicas e intercepto. Mostrar solo los términos de mayor magnitud es una explicación parcial: la suma completa debe incluir los restantes. Por ello, el material conserva tanto el modelo representativo como las predicciones, en lugar de ofrecer únicamente una descripción narrativa del resultado.
